\section[Stage 019]{Stage 019: Dimensionless Support-Feasibility Frontier}
\label{stage:019}

\stagefield{Part / anchor}{Part II; primary anchor \resultanchor{MTDC-T5}.}
\claimstatus{The support frontier is \StatusExactClosure{} inside the D/N selected-branch normal form.  Actual physical placement on the admissible frontier is \StatusOpen{}.}

\stagefield{Purpose}{Stage~019 isolates the support-feasibility condition that remains after the normalization locus has selected \(\xi_{\rm req}\).  The total loading must exceed the mixed-sector baseline by a nonnegative support/BdG contribution.  This requirement becomes the exact frontier \(M_{\rm mix}\leq G(\xi_{\rm req},\delta)\).}

\stagefield{Inputs}{The inputs are the Stage~018 total loading \(\alpha_{\rm req}(\xi,
\delta)\), the mixed-sector baseline \(\alpha_{\rm mix}=\mathcal X^2/(\Omega_U^2\Delta_0)\), and the dimensionless variables \(\xi,\delta\).}

\stagefield{Verification}{SymPy audit: \StageFile{scripts/moving_throat_pde_stage019_support_feasibility_frontier_sympy_audit.py}; transcript: \StageFile{scripts/output/moving_throat_pde_stage019_support_feasibility_frontier_sympy_audit.txt}.  Mathematica audit: \StageFile{mathematica/moving_throat_pde_stage019_support_feasibility_frontier_mathematica_audit.wl}; transcript: \StageFile{mathematica/output/moving_throat_pde_stage019_support_feasibility_frontier_mathematica_audit.txt}.}

\subsection*{Derivation}

\paragraph{1. Dimensionless mixed baseline and support function.}
Split
\begin{equation}
\alpha_{\rm req}=\alpha_{\rm mix}+\frac{g_{B,{\rm req}}^2}{\varpi^2},
\qquad
\alpha_{\rm mix}=\frac{\mathcal X^2}{\Omega_U^2\Delta_0}.
\end{equation}
Define
\begin{equation}
\label{eq:app-stage019-Mmix}
\boxed{
M_{\rm mix}:=\frac{8\alpha_{\rm mix}}{\pi^2A}
=
\frac{8\mathcal X^2}{\pi^2A\Omega_U^2\Delta_0}}.
\end{equation}
The dimensionless support-feasibility function is
\begin{equation}
\label{eq:app-stage019-G}
\boxed{
G(\xi,\delta):=\frac{8\alpha_{\rm req}}{\pi^2A}
=
\frac{9\xi(\xi+\delta)}{9\delta+11\xi}}.
\end{equation}
Therefore
\begin{equation}
\label{eq:app-stage019-gBreq}
\boxed{
\frac{g_{B,{\rm req}}^2}{\varpi^2}
=\frac{\pi^2A}{8}\left[G(\xi_{\rm req},\delta)-M_{\rm mix}\right]}.
\end{equation}
The support/BdG sector is feasible exactly when
\begin{equation}
\label{eq:app-stage019-feasible}
\boxed{
M_{\rm mix}\leq G(\xi_{\rm req},\delta)}.
\end{equation}

\paragraph{2. Monotonicity and endpoints.}
The derivative is
\begin{equation}
\label{eq:app-stage019-G-derivative}
\boxed{
\frac{\partial G}{\partial\xi}
=
\frac{9(9\delta^2+18\delta\xi+11\xi^2)}{(9\delta+11\xi)^2}>0}.
\end{equation}
The endpoints are
\begin{equation}
\label{eq:app-stage019-G-endpoints}
G(0,\delta)=0,
\qquad
G_{\max}(\delta)=\lim_{\xi\to1^-}G(\xi,\delta)
=\frac{9(1+\delta)}{9\delta+11}.
\end{equation}
The maximum is the dimensionless form of the refined stability bound on the mixed-sector baseline.

\paragraph{3. Parametric admissibility frontier.}
Stages~018 and 019 together define the parametric stable frontier
\begin{equation}
\label{eq:app-stage019-param-frontier}
\boxed{
\xi\in[0,1)
\longmapsto
\left(F(\xi,\delta),G(\xi,\delta)\right)}.
\end{equation}
For fixed \(\delta\), the normalization target chooses \(\xi_{\rm req}\) by
\begin{equation}
\label{eq:app-stage019-Fselect}
F(\xi_{\rm req},\delta)=R_{\rm target},
\end{equation}
and the support baseline must lie below the frontier at that same point.

Near onset,
\begin{equation}
\label{eq:app-stage019-G-onset-expansion}
G(\xi,
\delta)=\xi-\frac{2\xi^2}{9\delta}+O(\xi^3).
\end{equation}
Combined with \eqref{eq:app-stage018-xi-near-onset},
\begin{equation}
M_{\rm crit}\simeq\frac{R_{\rm target}-1}{1+8/(9\delta)}
\end{equation}
for targets just above onset.

\paragraph{4. Final admissibility test.}
The final Stage~019 branch test is
\begin{equation}
\label{eq:app-stage019-final-test}
\boxed{
R_{\rm target}\geq1,
\qquad
F(\xi_{\rm req},\delta)=R_{\rm target},
\qquad
M_{\rm mix}\leq G(\xi_{\rm req},\delta)}.
\end{equation}
If \(R_{\rm target}<1\), the branch cannot hit the target on the stable side.  If \(R_{\rm target}\geq1\) but \(M_{\rm mix}>G(\xi_{\rm req},\delta)\), the required support loading is negative and the physical branch is inadmissible in this reduced closure.

\stagefield{Output}{Stage~019 outputs the dimensionless mixed baseline \eqref{eq:app-stage019-Mmix}, the support-feasibility function \eqref{eq:app-stage019-G}, the required support loading \eqref{eq:app-stage019-gBreq}, the monotonicity derivative \eqref{eq:app-stage019-G-derivative}, and the final admissibility test \eqref{eq:app-stage019-final-test}.}

\stagefield{Checks}{\begin{verificationchecklist}
\item Substituting \(\alpha_{\rm req}\) from Stage~018 into \(8\alpha_{\rm req}/(\pi^2A)\) gives \(G\).
\item \(\partial_\xi G>0\) for \(\delta>0\) and \(0\leq\xi<1\).
\item \(M_{\rm mix}\leq G\) is exactly the nonnegativity of \(g_{B,{\rm req}}^2\).
\item The parametric map \((F,G)\) is one-dimensional at fixed \(\delta\), so the selected reduced branch has no remaining free hidden eigenvector parameter at this order.
\end{verificationchecklist}}

\stagefield{Downstream use}{Part~III begins by computing continuum-kernel placement data that feed \(A\), \(\Delta K_{\rm ax}\), \(\beta_0\), \(M_{\rm mix}\), and hence \(R_{\rm target}\), \(\delta\), and the exact frontier.  Later support-completion stages ask whether rank-two support and coherent D/N support can place the actual branch inside the admissible region.}

